\section{A fixed reflected boundary representation}\label{sec:boundary}
\subsection{Finite slab and reference color fiber}
For general SU($N_c$) on a finite slab, the Wilson action \cite{wilson} is
\begin{equation}
S=\beta\sum_p\left(1-\frac1{N_c}\Ree\tr U_p\right).
\end{equation}
Use normalized Haar measure on each link. Let U denote the central-slice
spatial links and $X_\pm$ the strict half-slab variables. Site reflection
in physical time gives
\begin{equation}
S(U,X_+,X_-)=S_0(U)+S_+(U,X_+)+S_-(U,X_-),
\qquad S_-=\theta S_+.
\end{equation}
There is no plaquette containing both strict half interiors. Define
\begin{equation}\label{eq:boundary-weight}
h(U)=\int e^{-S_+(U,X_+)}\dd X_+,
\quad \Omega(U)=e^{-S_0(U)/2}h(U),
\quad \dd\nu=Z^{-1}\Omega^2\dd m=w\dd m.
\end{equation}
The compact finite integral shows that $\Omega$ is smooth, gauge
invariant and strictly positive, with positive lower and upper bounds.
This is a prepared finite-slab amplitude with unit outer-boundary data.
It is consistent with the standard lattice transfer-matrix framework
\cite{luscher}; the gluing identity needed here is derived directly below.

At a fixed base vertex, retain a reference endomorphism fiber, so a
matrix-valued observable transforms by $F\mapsto gFg^{-1}$ there.
For a half-slab observable F set
\begin{equation}
\psi_F(U)=e^{-S_0(U)/2}\int e^{-S_+(U,X_+)}F(U,X_+)\dd X_+.
\end{equation}
Reflected gluing and contraction of the reference indices give
\begin{equation}\label{eq:gluing}
\ip F G_{\rm refl}=Z^{-1}\int\frac1{N_c}
             \tr(\psi_F^\dagger\psi_G)\dd m.
\end{equation}
For boundary functions B,C this reduces exactly to
\begin{equation}\label{eq:physical-H}
\ip B C_\nu=\int\frac1{N_c}\tr(B^\dagger C)\dd\nu,
\qquad\cH=L^2_{\rm cov}(M,\nu;\operatorname{End}\C^{N_c}).
\end{equation}
The reference fiber retains singlet and adjoint channels and adds no
bulk dynamical matter. Tracing each observable before taking its norm
would be a different scalar reduction.

\begin{proposition}[Boundary descent through the reflected null space]
The reflected null space is the kernel of $F\mapsto\psi_F$.
For every bounded covariant boundary multiplier D,
$\psi_{DF}=D\psi_F$. It therefore descends to the reflected quotient;
if $D^\dagger D=I$, its action is unitary. Based spatial loops Q have
$\norm Q_\nu=1$, and their Gram entries are
$\mathbb E_\nu\tr(Q_i^\dagger Q_j)/N_c$.
\end{proposition}
\begin{proof}
Separate the two half integrals in \eqref{eq:gluing}. Boundary
multiplication passes through the half integral, and the remaining
statements follow from the positive integral norm. Strict positivity
of $\Omega$ identifies its completed covariant image with
\eqref{eq:physical-H}.
\end{proof}
An interior multiplier need not descend: take independent signs $X_+$
and $X_-$ with reflected pairing
$\mathbb E\overline{F(X_-)}G(X_+)=\overline{\mathbb EF}\mathbb EG$.
The function $F(X_+)=X_+$ is null, while multiplication by the unitary
sign $X_+$ sends it to 1. Boundary dependence is the hypothesis that
prevents this failure. Slice-only gauge-covariant group-valued smoothing
preserves boundary admissibility; its bias and continuum limit still
require analysis. General four-dimensional smoothing does not
implicitly preserve the same reflection support.

\subsection{Weighted derivatives and compact source norms}
Return to SU(2). Normalize its bi-invariant metric so
$\ii\sigma_1,\ii\sigma_2,\ii\sigma_3$ are orthonormal; the fundamental
Casimir is 3. For Haar-divergence-free link derivatives $D_{\ell,j}$,
\begin{equation}\label{eq:electric}
D_{\ell,j}^{*\nu}=-D_{\ell,j}-2D_{\ell,j}\log\Omega,
\quad \mathfrak e_\nu(B,C)=\sum_{\ell,j}\ip{D_{\ell,j}B}{D_{\ell,j}C}_\nu,
\quad E_\nu=\sum_{\ell,j}D_{\ell,j}^{*\nu}D_{\ell,j}.
\end{equation}
Take the closed form and its nonnegative Friedrichs operator. The full
summed form is gauge invariant, so the covariant subspace reduces it;
an individual colored derivative need not preserve that subspace.
Ellipticity on compact M implies compact resolvent. For $r>0$,
\begin{equation}\label{eq:compact-tail}
\cK_r=\Dom(1+E_\nu)^{r/2}\Subset\cH,
\quad
\norm{(1-\mathbf1_{[0,L]}(E_\nu))B}
\le(1+L)^{-r/2}\norm{(1+E_\nu)^{r/2}B}.
\end{equation}
The inequality follows eigenvalue by eigenvalue and proves compactness
by finite-rank approximation. These are physical regularity norms, not
arithmetic translation generators.

\subsection{The one-plaquette marginal}
Fix an elementary spatial plaquette P using $\ell=4$ distinct links.
Its product map onto SU(2) is a smooth submersion: solve for one link
with the other three fixed. Disintegration of the smooth positive
weight, and differentiation under the compact fiber integral, give
\begin{equation}\label{eq:marginal}
P_*\nu=\rho(g)\dd g,\qquad
0<c_\rho\le\rho\le C_\rho<\infty.
\end{equation}
Gauge invariance makes $\rho$ central. Its radial extension is even,
smooth and $2\pi$-periodic. The class sector $F(P)I$ is a closed
subspace $\cH_{\rm cl}$ of \eqref{eq:physical-H}. Put
\begin{equation}
\chi_n(\theta)=\frac{\sin n\theta}{\sin\theta},\quad n\ge1,
\qquad t_k=\frac1\pi\int_0^\pi\rho(\theta)\cos(k\theta)\dd\theta.
\end{equation}
Then $t_{-k}=t_k\in\R$ decreases rapidly, $t_0-t_2=1$, and
\begin{equation}\label{eq:gram}
\ip{\chi_n}{\chi_m}_\nu=t_{n-m}-t_{n+m},\qquad
c_\rho\sum|c_n|^2\le\norm{\sum c_n\chi_n}_\nu^2
                       \le C_\rho\sum|c_n|^2.
\end{equation}
These statements follow from the radial Haar measure
$(2/\pi)\sin^2\theta\dd\theta$ and the product-to-sum identity.
In general $t_0\ne1$. The full $E_\nu$ need not preserve
$\cH_{\rm cl}$; its form compression must be distinguished from a
restriction to a reducing subspace.

\subsection{What full finite Ward equations determine}
For the full finite lattice, suppose a normalized positive measure
$\mu$ satisfies
$\int DF\dd\mu=\int F(DS)\dd\mu$ for every smooth F and every
link derivative D. Then $\eta=e^S\mu$ satisfies $\int Dh\dd\eta=0$,
by applying the identity to $F=e^S h$. Hence $\eta$ is invariant under
all link translations and is a constant multiple of product Haar.
Thus $\mu=Z^{-1}e^{-S}m$, without a prior density assumption.
This determines the state, not an arithmetic source law. A finite list
of loop equations does not have the stated all-test hypothesis.
